% !TEX root = ../main.tex
\section{Pseudocódigo de las funciones principales}

El pseudocódigo se escribe en sintaxis C simplificada: condicionales,
ciclos y llamadas a función explícitas, sin declarar tipos de bajo nivel
que no aportan a entender el algoritmo. Cada bloque corresponde a un solo
camino de decisión por línea, para que sea directo de traducir a un
diagrama de flujo (ver sección de diagramas): cada \texttt{si}, cada
ciclo y cada \texttt{retornar} es un bloque separado que corresponde a un
nodo o una decisión distinta en el diagrama.

% El pseudocodigo de cada funcion vive en su propio archivo dentro de
% informe/pseudocodigo/, editable sin tocar la prosa de esta seccion.

\subsection{registrarEstudiante(datos)}
\begin{lstlisting}[style=pseudocodigo]
funcion registrarEstudiante(datos):
    si no validarFormato(datos):
        retornar ERROR

    posicion = buscarEspacioLibre()
    si posicion == NO_HAY_ESPACIO:
        posicion = expandirVector()

    escribirRegistroEnOffset(identidad.dat, posicion, datos.identidad)
    escribirRegistroEnOffset(academico.dat, posicion, datos.academico)
    escribirRegistroEnOffset(contacto.dat, posicion, datos.contacto)

    insertarEnIndicePrimario(datos.carne, posicion)

    retornar EXITO
\end{lstlisting}


\subsection{buscarEstudiantePorCarne(carne)}
\begin{lstlisting}[style=pseudocodigo]
funcion buscarEstudiantePorCarne(carne):
    posicion = buscarPosicionPorCarne(carne)

    si posicion == NO_ENCONTRADO:
        retornar NO_ENCONTRADO

    identidad = leerRegistroEnOffset(identidad.dat, posicion)
    academico = leerRegistroEnOffset(academico.dat, posicion)
    contacto  = leerRegistroEnOffset(contacto.dat, posicion)

    retornar ensamblarEstudiante(identidad, academico, contacto)
\end{lstlisting}


\subsection{buscarPosicionPorCarne(carne)}
\begin{lstlisting}[style=pseudocodigo]
funcion buscarPosicionPorCarne(carne):
    posicion = busquedaBinaria(zonaOrdenada, carne)

    si posicion != NO_ENCONTRADO:
        retornar posicion

    para i desde 0 hasta longitud(listaRecienIngresados) - 1:
        si listaRecienIngresados[i].carne == carne:
            retornar listaRecienIngresados[i].posicion

    retornar NO_ENCONTRADO
\end{lstlisting}


\subsection{buscarEstudiantesPorApellido(apellido)}
\begin{lstlisting}[style=pseudocodigo]
funcion buscarEstudiantesPorApellido(apellido):
    si indiceApellido.desactualizado:
        indiceApellido = construirIndiceSecundario(APELLIDO)

    posiciones = buscarEnIndiceSecundario(indiceApellido, apellido)

    resultados = listaVacia()
    para cada posicion en posiciones:
        identidad = leerRegistroEnOffset(identidad.dat, posicion)
        agregar(resultados, identidad)

    retornar resultados
\end{lstlisting}


\subsection{eliminarEstudianteLogico(carne)}
\begin{lstlisting}[style=pseudocodigo]
funcion eliminarEstudianteLogico(carne):
    posicion = buscarEnIndicePrimario(carne)

    si posicion == NO_ENCONTRADO:
        retornar ERROR

    marcarInactivoEnIndice(carne)

    academico = leerRegistroEnOffset(academico.dat, posicion)
    academico.estado = 0
    escribirRegistroEnOffset(academico.dat, posicion, academico)

    retornar EXITO
\end{lstlisting}


\subsection{registrarPosicionEnMapa(carne, posicion)}
\begin{lstlisting}[style=pseudocodigo]
funcion insertarEnIndicePrimario(carne, posicion):
    agregarAlFinal(zonaAltasNuevas, carne, posicion)

    si longitud(zonaAltasNuevas) >= 0.10 * longitud(zonaPrimaria):
        reorganizarIndicePrimario()
\end{lstlisting}


\subsection{fusionarRecienIngresados()}
\begin{lstlisting}[style=pseudocodigo]
funcion reorganizarIndicePrimario():
    fusionado = mezclarOrdenado(zonaPrimaria, zonaAltasNuevas)
    zonaPrimaria = fusionado
    zonaAltasNuevas = vacio()
\end{lstlisting}


\subsection{purgarInactivos()}
\begin{lstlisting}[style=pseudocodigo]
funcion purgarInactivos():
    crear identidadTemp.dat, academicoTemp.dat, contactoTemp.dat

    para cada posicion desde 0 hasta contarRegistros() - 1:
        academico = leerRegistroEnOffset(academico.dat, posicion)

        si academico.marcadoParaPurga != 1:
            identidad = leerRegistroEnOffset(identidad.dat, posicion)
            contacto  = leerRegistroEnOffset(contacto.dat, posicion)
            agregarRegistroAlFinal(identidadTemp.dat, identidad)
            agregarRegistroAlFinal(academicoTemp.dat, academico)
            agregarRegistroAlFinal(contactoTemp.dat, contacto)
        // si marcadoParaPurga es 1, el registro se descarta y no pasa al archivo nuevo
        // (independiente de estadoAcademico: un estudiante en pausa o con permiso
        // no se marca para purga, asi que nunca se elimina solo por su estado academico)

    reemplazarArchivo(identidad.dat, identidadTemp.dat)
    reemplazarArchivo(academico.dat, academicoTemp.dat)
    reemplazarArchivo(contacto.dat, contactoTemp.dat)

    // las posiciones cambiaron: reconstruir los tres mapas desde cero,
    // no se puede reutilizar la lista de recien ingresados anterior
    mapaCarne = reconstruirMapa(identidad.dat, CARNE)
    mapaIdentificacion = reconstruirMapa(identidad.dat, IDENTIFICACION)
    mapaApellido = reconstruirMapa(identidad.dat, APELLIDO)

    retornar EXITO
\end{lstlisting}

